\section{Conclusion}
\label{sec:conclusion}

This paper introduced \tool, a framework for checking protocol implementations
against natural-language standards through field-centric reasoning and
test-based validation. The central idea is to use protocol fields as the
semantic bridge between prose and code: standards are reduced to structured
\ruleform{} rules, rules are aligned to implementation variables and contexts,
and consistency judgments are made only after the evidence is sufficient or the
remaining uncertainty is explicit. By combining constrained LLM reasoning,
syntax-aware program analysis, and executable validation tests, \tool{} aims to
make LLM-assisted conformance checking auditable and reproducible. Our
evaluation produced 204 conformance-deviation reports across widely used
protocol implementations. Maintainers have confirmed and fixed 126 of these
reports, roughly nine were false positives, and 69 are still awaiting response
or final adjudication; two confirmed issues became previously unknown CVEs. The
average GPT-5.4 model cost is \$1.122 per submitted report, indicating that
constrained LLM reasoning can be practical when paired with field-level
specifications and reproducible validation artifacts.
